\begin{table}[H]
\centering
\caption{Nama Daerah Matriks Transisi Kategori IPM (UNDP) 2021 \(\rightarrow\) 2022} \label{tab:sample2ipmname2021x2022}
\begin{adjustbox}{max width=\textwidth}
\begin{tabular}{|p{4cm}|p{10cm}|}
\hline
\textbf{Kategori} & \textbf{Daftar Wilayah} \\
\hline
Sangat Tinggi (2021) \(\rightarrow\) Sangat Tinggi (2022) & Kota Jakarta Pusat \\
\hline
Sangat Tinggi (2021) \(\rightarrow\) Tinggi (2022) & - \\
\hline
Sangat Tinggi (2021) \(\rightarrow\) Sedang (2022) & - \\
\hline
Sangat Tinggi (2021) \(\rightarrow\) Rendah (2022) & - \\
\hline
Tinggi (2021) \(\rightarrow\) Sangat Tinggi (2022) & - \\
\hline
Tinggi (2021) \(\rightarrow\) Tinggi (2022) & Kota Tasikmalaya, Kota Cilegon \\
\hline
Tinggi (2021) \(\rightarrow\) Sedang (2022) & - \\
\hline
Tinggi (2021) \(\rightarrow\) Rendah (2022) & - \\
\hline
Sedang (2021) \(\rightarrow\) Sangat Tinggi (2022) & - \\
\hline
Sedang (2021) \(\rightarrow\) Tinggi (2022) & - \\
\hline
Sedang (2021) \(\rightarrow\) Sedang (2022) & Bandung Barat \\
\hline
Sedang (2021) \(\rightarrow\) Rendah (2022) & - \\
\hline
Rendah (2021) \(\rightarrow\) Sangat Tinggi (2022) & - \\
\hline
Rendah (2021) \(\rightarrow\) Tinggi (2022) & - \\
\hline
Rendah (2021) \(\rightarrow\) Sedang (2022) & - \\
\hline
Rendah (2021) \(\rightarrow\) Rendah (2022) & - \\
\hline
\end{tabular}
\end{adjustbox}
\end{table}